(1) Conjugation by the invertible element $e_+$ is an algebra automorphism of $\Cl(p{+}1,q{+}1)$ (on vectors, as $e_+^{2} = 1$, it is minus the reflection in $e_+^{\perp}$, \claimref{std-algebra.reflection}); in particular it is injective, so it sends nonzero elements to nonzero elements. By \eqref{eq:action} and linearity it maps vectors of the ambient space to vectors, and it preserves squared lengths: $(e_+ N e_+)^2 = e_+ N^2 e_+ = N^2$ for any vector $N$, since $N^2$ is a scalar. Hence it maps nonzero null vectors to nonzero null vectors, and rays to rays ($\lambda N \mapsto \lambda\, e_+Ne_+$), so $\iota$ is well defined and total on $\mathcal{Q}$. It is continuous because it descends from an invertible linear map of $\mathbb{V} \oplus H$ through the quotient $\pi$. Finally $e_+ (e_+ N e_+) e_+ = N$ since $e_+^2 = 1$, so $\iota^2 = \mathrm{id}$.

(2) By linearity of conjugation and \eqref{eq:action},
\begin{equation}\label{eq:keydisplay}
e_+ X(x)\, e_+ \;=\; -\Bigl( Q(x)\,\nO + x + \tfrac12 \nInf \Bigr) \qquad \text{for every } x \in \mathbb{V}.
\end{equation}
For $Q(x) \neq 0$, rescaling by $-1/Q(x)$ gives $\nO + x/Q(x) + \tfrac12 (1/Q(x)) \nInf$, and since $Q(x/Q(x)) = 1/Q(x)$ this is exactly $X(x/Q(x))$ by \eqref{eq:lift}.

(3) For $Q(x_0) = 0$ the display \eqref{eq:keydisplay} gives $e_+ X(x_0) e_+ = -(x_0 + \tfrac12 \nInf)$, a nonzero vector (its $\nInf$-coefficient is $\tfrac12$) whose $\nO$-coefficient vanishes: an ideal point, null since $x_0 \cdot \nInf = 0$ and $x_0^2 = \nInf^2 = 0$; for $x_0 = 0$ it is the class $[\nInf]$. For the limit claim, consider the map $J : \mathbb{V} \to \mathcal{Q}$, $J(x) = [\,Q(x)\nO + x + \tfrac12\nInf\,]$. Its representative is polynomial in $x$ and never zero (the $\nInf$-coefficient is $\tfrac12$), so $J$ is continuous; and by \eqref{eq:keydisplay}, $J(x) = \iota[X(x)]$ for \emph{every} $x \in \mathbb{V}$, whether $Q(x)$ vanishes or not. Continuity of $J$ at $x_0$ now gives $\iota[X(x)] \to \iota[X(x_0)]$ along every path of invertible $x \to x_0$, as claimed. For uniqueness: the lifted invertible points are dense in $\mathcal{Q}$ (Lemma \claimref{transport.density}) and $\mathcal{Q}$ is Hausdorff, so any two continuous maps $\mathcal{Q} \to \mathcal{Q}$ that agree on them agree everywhere; $\iota$ agrees with the lifted reciprocal there by (2).
